%============================================================
% CH1 · P5 · PROBLEM 2 — UNIQUE NUMBER OF OCCURRENCES (LeetCode 1207)
%============================================================
\subsection{Problem 1.5.2: Unique Number of Occurrences}
\problemlabel{1.5.2}{ssec:ch1-unique-occurrences}

\problemcard{1.5.2}{Unique Number of Occurrences}{LeetCode 1207}{https://leetcode.com/problems/unique-number-of-occurrences/}{\exeasy\ \lc{Easy}}{Counting array with an offset, then a frequency array of the counts (skip zero counts)}

\subsubsection*{1. Problem statement}

Given an integer array \texttt{arr}, return \texttt{true} if the number of occurrences of each distinct value is different from the number of occurrences of every other distinct value, otherwise \texttt{false}.

\textbf{Constraints:} $1 \le n \le 1000$, \; $-1000 \le \texttt{arr[i]} \le 1000$.

\begin{table}[H]
\centering\small
\begin{tabular}{@{}lll@{}}
\toprule
\textbf{arr} & \textbf{Output} & \textbf{Counts} \\
\midrule
$[1,2,2,1,1,3]$ & \texttt{true} & $1{:}\,3$, \ $2{:}\,2$, \ $3{:}\,1$ --- all different \\
$[1,2]$ & \texttt{false} & $1{:}\,1$, \ $2{:}\,1$ --- the count 1 appears twice \\
$[-3,0,1,-3,1,1,1,-3,10,0]$ & \texttt{true} & $-3{:}\,3$, \ $0{:}\,2$, \ $1{:}\,4$, \ $10{:}\,1$ \\
\bottomrule
\end{tabular}
\end{table}

\begin{diagrambox}
\begin{tikzpicture}[x=1cm, y=1cm]
    \dparray{a}{0.9}{arr}{1/gvA, 2/gvB, 2/gvB, 1/gvA, 1/gvA, 3/acell}
    \node[note, anchor=west, text width=6.0cm, align=left] at (5.8,0.9) {Sample 1: value 1 (blue) three times, value 2 (red) twice, value 3 once};
    \begin{scope}[yshift=-1.0cm]
        \foreach \v/\c/\s [count=\k from 0] in {1/3/gvA, 2/2/gvB, 3/1/acell}{
            \node[\s, minimum size=7mm] (val\k) at (0.45+\k*1.75, 0) {\v};
            \node[anscell, minimum width=8mm] (cnt\k) at (1.6+\k*1.75, 0) {\c};
            \draw[flow] (val\k) -- (cnt\k);
        }
        \node[note, anchor=west, text width=6.0cm, align=left] at (5.8,0) {counts $3, 2, 1$: no count repeats $\Rightarrow$ \texttt{true}};
    \end{scope}
\end{tikzpicture}
\figcap{The question is about the \emph{counts}, not about the values: list one count per distinct value and check that the list has no repetition.}
\end{diagrambox}

\subsubsection*{2. The problem model}
\begin{modelbox}[Problem Model --- Unique Number of Occurrences]
\renewcommand{\arraystretch}{1.35}
\begin{tabularx}{\linewidth}{@{}>{\bfseries}l X@{}}
Input & an array of $n$ integers in $[-1000, 1000]$ \\
Computed & yes/no: is the map ``distinct value $\mapsto$ its count'' injective? \\
Objects & the distinct values ($\le \min(n, 2001)$) and their counts (integers in $[1, n]$) \\
Size & $n \le 1000$, value range $V = 2001$: two counting arrays of size $\bigO{V}$ and $\bigO{n}$, linear time \\
\end{tabularx}
\end{modelbox}

\begin{diagrambox}
\begin{tikzpicture}[x=1cm, y=1cm]
    \node[tag=primaryblue, text width=1.8cm] (A) at (0,0) {array\\$n \le 1000$};
    \node[tag=memoorange, text width=2.4cm] (C) at (3.4,0) {\texttt{cnt[v + 1000]}\\one per value};
    \node[tag=successgreen, text width=2.4cm] (S) at (7.3,0) {\texttt{seen[c]}\\one per count};
    \node[tag=accentorange, text width=1.8cm] (R) at (10.6,0) {any \texttt{seen[c]} $\ge 2$?};
    \draw[flow] (A) -- node[lbl, above]{count} (C);
    \draw[flow] (C) -- node[lbl, above]{skip 0s} (S);
    \draw[flow] (S) -- (R);
\end{tikzpicture}
\figcap{Two counting passes: values $\to$ counts, then counts $\to$ how many values have each count.}
\end{diagrambox}

\subsubsection*{3. Key observations}

\begin{itemize}[leftmargin=*]
    \item \textbf{Values can be negative:} store value $v$ at index $v + 1000$ (Tool~2 of the theory); 2001 cells.
    \item \textbf{Counts are small too:} every count lies in $[1, n]$, so ``do two values share a count?'' is a duplicate test on small integers --- another counting array, \texttt{seen[c]} (Tool~4: frequency of frequencies).
    \item \textbf{Tempting wrong idea --- scan every cell of \texttt{cnt}.} Most of the 2001 cells are 0 (values that never occur). Treating 0 as a count makes \emph{every} array with at least two absent values look like ``two values with count 0'' $\Rightarrow$ \texttt{false}. On Sample 1 this returns \texttt{false} instead of \texttt{true}. Only values that occur have an ``occurrence count''.
    \item \textbf{Equivalent formulation:} the answer is \texttt{true} iff (number of distinct values) $=$ (number of distinct counts). A set of counts compared with the number of non-zero cells gives the same answer.
\end{itemize}

\begin{diagrambox}
\begin{tikzpicture}[x=1cm, y=1cm]
    \dpindex{vv}{0.65}{value}{-1}{4}
    \dparray{c}{0}{\texttt{cnt}}{0/aout, 0/aout, 3/compcell, 2/compcell, 1/compcell, 0/aout}
    \node[note, anchor=west, text width=5.6cm, align=left] at (5.8,0.1) {grey cells: values that do not occur. Counting them as ``count 0'' finds a fake repetition ($0$ three times here, $1998$ times over the full range) and wrongly returns \texttt{false}.};
\end{tikzpicture}
\figcap{Counter-example to scanning all cells (part of the value range of Sample 1).}
\end{diagrambox}

\subsubsection*{4. Approaches}

\begin{bruteforcebox}
\textbf{Idea.} For every pair of positions holding different values, count both values by scanning the array; if the two counts are equal, answer \texttt{false}.
\begin{lstlisting}[style=cpp]
#include <bits/stdc++.h>
using namespace std;

class Solution {
public:
    bool uniqueOccurrences(vector<int>& arr) {
        int n = arr.size();
        for (int i = 0; i < n; i++)
            for (int j = 0; j < n; j++) {
                if (arr[i] == arr[j]) continue;    // need two different values
                int ci = count(arr.begin(), arr.end(), arr[i]);
                int cj = count(arr.begin(), arr.end(), arr[j]);
                if (ci == cj) return false;
            }
        return true;
    }
};

int main() {                                       // input: n, then the n values
    int n;
    if (!(cin >> n)) return 0;
    vector<int> a(n);
    for (auto& x : a) cin >> x;
    cout << (Solution().uniqueOccurrences(a) ? "true" : "false") << "\n";
}
\end{lstlisting}
\begin{complexitybox}
\textbf{Time:} $\bigO{n^3}$ --- $n^2$ pairs, each with two $\bigO{n}$ scans: up to $2\cdot10^9$ steps for $n = 1000$, \tle{} on large tests with many distinct values. \textbf{Space:} $\bigO{1}$. It recounts the same values again and again --- exactly what a counting array stores once.
\end{complexitybox}
\end{bruteforcebox}

\begin{algorithmbox}
\textbf{Idea.} Count the values with an offset. Then scan the counting array; for every non-zero count $c$, increment \texttt{seen[c]}; a second value with the same count makes \texttt{seen[c]} reach 2.
\begin{lstlisting}[style=cpp]
#include <bits/stdc++.h>
using namespace std;

class Solution {
public:
    bool uniqueOccurrences(vector<int>& arr) {
        const int OFF = 1000;                      // values -1000..1000 -> indices 0..2000
        vector<int> cnt(2 * OFF + 1, 0);
        for (int x : arr) cnt[x + OFF]++;
        vector<int> seen(arr.size() + 1, 0);       // seen[c] = how many values occur c times
        for (int c : cnt) {
            if (c == 0) continue;                  // absent values are not "occurrences"
            if (++seen[c] > 1) return false;       // a second value with the same count
        }
        return true;
    }
};

int main() {                                       // input: n, then the n values
    int n;
    if (!(cin >> n)) return 0;
    vector<int> a(n);
    for (auto& x : a) cin >> x;
    cout << (Solution().uniqueOccurrences(a) ? "true" : "false") << "\n";
}
\end{lstlisting}
\begin{complexitybox}
\textbf{Time:} $\bigO{n + V}$ with $V = 2001$. \textbf{Space:} $\bigO{V + n}$. \textbf{Verdict:} \ac{}.
\end{complexitybox}
\end{algorithmbox}

\begin{invariantbox}[Why it is correct]
After the first loop, $\texttt{cnt}[v + 1000]$ is the number of occurrences of $v$ (Tool~1 with an offset), and the distinct values of the array are exactly the $v$ with a non-zero cell. During the second loop, $\texttt{seen}[c]$ equals the number of distinct values scanned so far whose count is $c$. If some count is shared by two distinct values, then when the second of them is scanned \texttt{seen[c]} becomes 2 and we return \texttt{false}; if the loop finishes, every \texttt{seen[c]} is at most 1, i.e.\ the counts are pairwise different. Indices are safe: $1 \le c \le n$ and \texttt{seen} has $n+1$ cells. \hfill$\blacksquare$
\end{invariantbox}

\begin{alternativebox}
\textbf{Sort, measure runs, sort the run lengths.} After sorting, equal values form consecutive runs, and the length of a run is the count of its value --- no value range needed. Collect the run lengths, sort them, and look for two equal neighbours.
\begin{lstlisting}[style=cpp]
#include <bits/stdc++.h>
using namespace std;

class Solution {
public:
    bool uniqueOccurrences(vector<int>& arr) {
        vector<int> a(arr);
        sort(a.begin(), a.end());                  // equal values become one run
        vector<int> runs;
        for (size_t i = 0, j; i < a.size(); i = j) {
            for (j = i; j < a.size() && a[j] == a[i]; j++) {}
            runs.push_back(j - i);                 // length of the run of a[i]
        }
        sort(runs.begin(), runs.end());
        for (size_t t = 1; t < runs.size(); t++)
            if (runs[t] == runs[t - 1]) return false;   // two runs of equal length
        return true;
    }
};

int main() {                                       // input: n, then the n values
    int n;
    if (!(cin >> n)) return 0;
    vector<int> a(n);
    for (auto& x : a) cin >> x;
    cout << (Solution().uniqueOccurrences(a) ? "true" : "false") << "\n";
}
\end{lstlisting}
$\bigO{n \log n}$ time, $\bigO{n}$ memory, and it works unchanged for values up to $10^9$ (variant V4). The run-walking loop (an inner pointer $j$ that jumps over a block of equal values) returns in Chapter~\ref{ch:two_pointers}.

\begin{diagrambox}
\begin{tikzpicture}[x=1cm, y=1cm]
    \dparray{s}{0}{sorted}{-3/awin, -3/awin, -3/awin, 0/akeep, 0/akeep, 1/aread, 1/aread, 1/aread, 1/aread, 10/acell}
    \arange{s0}{s2}{3}{memoorange}
    \arange{s3}{s4}{2}{successgreen}
    \arange{s5}{s8}{4}{accentorange}
    \arange{s9}{s9}{1}{darkgray}
\end{tikzpicture}
\figcap{Sample 3 sorted: run lengths $3, 2, 4, 1$; sorted $1, 2, 3, 4$ --- no equal neighbours $\Rightarrow$ \texttt{true}.}
\end{diagrambox}
\end{alternativebox}

\begin{tricksbox}
\begin{itemize}[leftmargin=*]
    \item \textbf{A boolean array is enough for \texttt{seen}:} \texttt{if (seen[c]) return false; seen[c] = true;}.
    \item \textbf{Pigeonhole shortcut:} $d$ distinct values with pairwise different counts need at least $1 + 2 + \dots + d = d(d+1)/2$ elements. With $n \le 1000$, more than 44 distinct values ($44 \cdot 45 / 2 = 990$, $45 \cdot 46/2 = 1035$) already force \texttt{false}.
    \item \textbf{\texttt{unordered\_map} + \texttt{unordered\_set}} is the one-size-fits-all version (Chapter~\ref{ch:hashing}); here the arrays are simpler and faster.
\end{itemize}
\end{tricksbox}

\subsubsection*{5. Trace}

\begin{tracebox}[{Sample 3: $[-3,0,1,-3,1,1,1,-3,10,0]$}]
\centering
\begin{tikzpicture}[x=1cm, y=1cm]
    \dparray{a}{2.3}{arr}{-3/acell, 0/acell, 1/acell, -3/acell, 1/acell, 1/acell, 1/acell, -3/acell, 10/acell, 0/acell}
    \node[font=\scriptsize\bfseries, text=deepblue] at (0.45,1.35) {cell};
    \foreach \v/\idx/\c [count=\k from 0] in {-3/997/3, 0/1000/2, 1/1001/4, 10/1010/1}{
        \node[idx] at (1.6+\k*1.9, 1.35) {\idx\ ($v = \v$)};
        \node[compcell] (c\k) at (1.6+\k*1.9, 0.75) {\c};
    }
    \node[rowlabel] at (-0.15,0.75) {\texttt{cnt}};
    \node[note] at (8.9,0.75) {all other cells 0};
    \dpindex{sv}{-1.55}{count $c$}{0}{5}
    \dparray{s}{-0.9}{\texttt{seen}}{0/emptycell, 1/basecell, 1/basecell, 1/basecell, 1/basecell, 0/emptycell}
    \draw[dep] (c0.south) to[bend left=10] (s3.north);
    \draw[dep] (c1.south) to[bend right=10] (s2.north);
    \draw[dep] (c2.south) to[bend left=10] (s4.north);
    \draw[dep] (c3.south) to[bend right=20] (s1.north);
\end{tikzpicture}
\figcap{Only four cells of \texttt{cnt} are non-zero. Each of them marks its own count in \texttt{seen}; no cell of \texttt{seen} reaches 2.}

\medskip
\small
\begin{tabular}{@{}cccccl@{}}
\toprule
\textbf{step} & \textbf{value $v$} & \textbf{cell $v + 1000$} & $c = \texttt{cnt}$ & \texttt{seen[c]} after & \textbf{verdict so far} \\
\midrule
1 & $-3$ & 997 & 3 & 1 & ok \\
2 & $0$ & 1000 & 2 & 1 & ok \\
3 & $1$ & 1001 & 4 & 1 & ok \\
4 & $10$ & 1010 & 1 & 1 & ok \\
--- & \multicolumn{4}{l}{all other 1997 cells are 0 and skipped} & \texttt{true} \\
\bottomrule
\end{tabular}
\end{tracebox}

\begin{tracebox}[{Sample 2: $[1,2]$}]
\centering
\begin{tikzpicture}[x=1cm, y=1cm]
    \node[idx] at (1.0,1.35) {1001 ($v = 1$)};
    \node[idx] at (3.0,1.35) {1002 ($v = 2$)};
    \node[compcell] (d0) at (1.0,0.75) {1};
    \node[compcell] (d1) at (3.0,0.75) {1};
    \node[rowlabel] at (-0.15,0.75) {\texttt{cnt}};
    \dpindex{tv}{-1.55}{count $c$}{0}{2}
    \dparray{t}{-0.9}{\texttt{seen}}{0/emptycell, 2/anscell, 0/emptycell}
    \draw[dep] (d0.south) to[bend right=15] (t1.north);
    \draw[dep] (d1.south) to[bend left=15] (t1.north);
    \node[note, anchor=west, text width=4.6cm, align=left] at (4.3,-0.1) {both values have count 1: the second one pushes $\texttt{seen}[1]$ to 2 $\Rightarrow$ \texttt{false} immediately};
\end{tikzpicture}

\medskip
\small
\begin{tabular}{@{}cccccl@{}}
\toprule
\textbf{step} & \textbf{value $v$} & \textbf{cell} & $c$ & \texttt{seen[c]} after & \textbf{verdict} \\
\midrule
1 & 1 & 1001 & 1 & 1 & ok \\
2 & 2 & 1002 & 1 & \textbf{2} & return \texttt{false} \\
\bottomrule
\end{tabular}
\end{tracebox}

\subsubsection*{6. Test suite}

\begin{testbox}
\centering\small
\begin{tabular}{@{}lll@{}}
\toprule
\textbf{arr} & \textbf{Why it is interesting} & \textbf{Expected} \\
\midrule
$[1,2,2,1,1,3]$ & official Sample 1 & \texttt{true} \\
$[1,2]$ & official Sample 2 & \texttt{false} \\
$[-3,0,1,-3,1,1,1,-3,10,0]$ & official Sample 3: negative values & \texttt{true} \\
$[5]$ & a single value & \texttt{true} \\
$[-1000, 1000, 1000]$ & both ends of the value range & \texttt{true} \\
$[4,4,7,7]$ & two values with count 2 & \texttt{false} \\
$[0,0,0]$ & value 0 (cell 1000) & \texttt{true} \\
\bottomrule
\end{tabular}

\medskip
The three programs agree on the official samples and on 500 random arrays ($n \le 10$, at most four distinct values chosen near $-1000$, $0$ and $1000$ to exercise the offset).
\end{testbox}

\subsubsection*{7. Common mistakes}

\begin{warningbox}[Mistake 1 --- no offset]
\texttt{cnt[x]++} with $x = -3$ writes outside the array. Use \texttt{cnt[x + 1000]} and a size of $2001$.
\end{warningbox}

\begin{warningbox}[Mistake 2 --- counting the zeros]
\begin{lstlisting}[style=snippet]
for (int c : cnt) if (++seen[c] > 1) return false;                 // WRONG: seen[0] explodes
for (int c : cnt) if (c > 0 && ++seen[c] > 1) return false;        // RIGHT
\end{lstlisting}
\end{warningbox}

\begin{warningbox}[Mistake 3 --- \texttt{seen} too small]
A count can be as large as $n$ (all elements equal). \texttt{seen} needs $n + 1$ cells, not 1001 by habit when $n$ may differ, and not $V$.
\end{warningbox}

\begin{warningbox}[Mistake 4 --- comparing counts of positions, not of values]
Looping over \emph{positions} $i < j$ and comparing $\texttt{cnt}[a_i]$ with $\texttt{cnt}[a_j]$ reports a clash for two copies of the same value ($[1,1]$ $\to$ \texttt{false}). The pairs must be pairs of \emph{distinct values}.
\end{warningbox}

\subsubsection*{8. Takeaways}

\begin{itemize}[leftmargin=*]
    \item Questions about multiplicities are answered by a counting array; questions about the multiplicities themselves by a second counting array over the counts.
    \item Negative values need an offset; absent values must be skipped.
    \item ``All different'' is a duplicate test --- on small integers, a boolean array decides it in one pass.
    \item Sorting turns counts into run lengths when the values are too large to index.
\end{itemize}

\begin{variantsbox}
\variant{Count Elements With Maximum Frequency}{\exeasy\ \lc{Easy}\ \varlink{Count Elements With Maximum Frequency}{LeetCode 3005}{https://leetcode.com/problems/count-elements-with-maximum-frequency/}}
  {return the total number of elements whose value has the maximum count (values $1..100$).}
  {model: $m = \max_v \texttt{cnt}[v]$; answer $m \cdot \texttt{fof}[m]$.}
  {algorithm: count, take the maximum, add the counts equal to it. Samples $\to 4, 5$.}
\variant{Sum of Unique Elements}{\exeasy\ \lc{Easy}\ \varlink{Sum of Unique Elements}{LeetCode 1748}{https://leetcode.com/problems/sum-of-unique-elements/}}
  {sum the values that occur exactly once.}
  {model: $\sum_v v \cdot [\texttt{cnt}[v] = 1]$.}
  {algorithm: count, then one pass over the 100 values. Samples $\to 4, 0, 15$.}
\variant{Most Frequent Even Element}{\exeasy\ \lc{Easy}\ \varlink{Most Frequent Even Element}{LeetCode 2404}{https://leetcode.com/problems/most-frequent-even-element/}}
  {the most frequent even value (smallest on ties), or $-1$; values up to $10^5$.}
  {model: a counting array restricted to even values; the mode with a tie rule.}
  {algorithm: count only even values, scan upwards keeping the first strict maximum. Samples $\to 2, 4, -1$.}
\variant{Large values}{\varnone}
  {the same question with $|a_i| \le 10^9$, $n \le 10^5$.}
  {model: counts still lie in $[1, n]$, but the values cannot index an array.}
  {algorithm: the alternative (sort, runs, sort the runs) in $\bigO{n \log n}$, or a hash map for the values plus the \texttt{seen} array for the counts.}
\end{variantsbox}
